% Lecture 1 notes: \input by notes.tex, \subimport'ed by the semester book.
% Exam questions (экзамен): 1 (DLP в Z/nZ, протокол Диффи–Хеллмана),
% 2 (тривиальная оценка сложности DLP, алгоритм Шенкса и его сложность).
\begin{lecture}{1}{Задача дискретного логарифмирования. Алгоритм Шенкса}

\section{Как договориться о ключе по открытому каналу}

Алиса и Боб хотят переписываться, используя быстрый симметричный шифр (например, AES), но для
этого им нужен общий секретный ключ. Встретиться лично они не могут, а весь их канал связи
прослушивает Ева. В 1976 г. У.~Диффи и М.~Хеллман~\cite{diffie-hellman} предложили способ
\emph{согласовать} общий секрет, обмениваясь только открытыми сообщениями. Протокол опирается
на асимметрию двух задач в группе $\Znx{p}$:
\begin{itemize}
  \item по $g$ и $x$ вычислить $g^x \bmod p$ легко "--- бинарным алгоритмом за $\bigO(\log x)$
    умножений (лекция~6 первого семестра);
  \item по $g$ и $h = g^x \bmod p$ найти показатель $x$ "--- \emph{дискретный логарифм} "---
    по-видимому, трудно: лучшие известные алгоритмы для $p$ длиной $2048$ бит практически
    неосуществимы.
\end{itemize}

Функция $x \mapsto g^x \bmod p$ "--- основной кандидат в \emph{односторонние функции}
(лекция~1 первого семестра). В этой лекции мы определим задачу дискретного логарифмирования
(DLP, discrete logarithm problem), разберём протоколы, стойкость которых на ней основана
(Диффи"--~Хеллман, Эль-Гамаль), и изучим два универсальных алгоритма её решения "---
тривиальный перебор за $\bigO(n)$ и алгоритм Шенкса за $\bigO(\sqrt n)$ операций, где $n$ "---
порядок основания. В следующих лекциях появятся алгоритмы, использующие структуру группы:
Полиг"--~Хеллман (лекция~2), $\rho$-метод Полларда (лекция~3), исчисление индексов (лекция~4).

\section{Напоминание: циклические группы и порядки}

Всюду $G$ "--- конечная группа с мультипликативной записью, $|G|$ "--- её порядок. Нам нужны
факты из лекции~7 первого семестра.

\begin{definition}
\emph{Порядок} элемента $g \in G$ "--- наименьшее $n \ge 1$, при котором $g^n = 1$; обозначение
$\ord(g)$. Подгруппа $\langle g \rangle = \{g^k : k \in \ZZ\}$ называется \emph{циклической
подгруппой}, порождённой $g$; группа $G$ \emph{циклическая}, если $G = \langle g \rangle$ для
некоторого $g$ (\emph{образующего}).
\end{definition}

\begin{proposition}\label{l1:order}
Пусть $\ord(g) = n$. Тогда
\begin{enumerate}
  \item $g^k = 1 \iff n \mid k$; в частности, $g^k = g^l \iff \congr{k}{l}{n}$;
  \item $\langle g \rangle = \{1, g, \dots, g^{n-1}\}$ и эти $n$ элементов различны;
  \item $n \mid |G|$ (теорема Лагранжа), поэтому $g^{|G|} = 1$.
\end{enumerate}
\end{proposition}

\begin{proof}
1. Разделим с остатком: $k = qn + r$, $0 \le r < n$. Тогда $g^k = (g^n)^q g^r = g^r$, и
$g^r = 1$ при $0 \le r < n$ только для $r = 0$ по минимальности $n$. Второе утверждение
следует из первого, применённого к $g^{k - l}$.
2. Любая степень $g^k$ равна $g^{k \bmod n}$, а $g^0, \dots, g^{n-1}$ различны по п.~1.
3. Это теорема Лагранжа для подгруппы $\langle g \rangle$, имеющей $n$ элементов: смежные классы
$a\langle g \rangle$ разбивают $G$ на части по $n$ элементов.
\end{proof}

\begin{theorem}[без доказательства, см. лекцию~7 первого семестра]\label{l1:cyclic}
Для простого $p$ группа $\Znx{p}$ циклическая порядка $p - 1$. Её образующие называются
\emph{первообразными корнями} по модулю $p$; их ровно $\varphi(p - 1)$.
\end{theorem}

Порядок $g \in \Znx{p}$ находится методом \emph{последовательного снижения}: начинаем
с $o = p - 1$ и для каждого простого $q \mid o$, пока $g^{o/q} = 1$, заменяем $o \gets o/q$.
Элемент $g$ "--- первообразный корень тогда и только тогда, когда $g^{(p-1)/q} \ne 1$ для всех
простых $q \mid p - 1$.

Следующий факт понадобится нам и в задачах, и при выборе параметров протоколов.

\begin{proposition}\label{l1:subgroups}
Пусть $G = \langle \gamma \rangle$ "--- циклическая группа порядка $N$ и $d \mid N$. Тогда
в $G$ есть ровно одна подгруппа порядка $d$:
\[
  H_d = \{y \in G : y^d = 1\} = \langle \gamma^{N/d} \rangle,
\]
и всякий элемент порядка $d$ её порождает. В частности, если $\ord(g) = n$, то $h \in \langle g
\rangle \iff h^n = 1$.
\end{proposition}

\begin{proof}
Пусть $y = \gamma^t$. По предложению~\ref{l1:order} $y^d = \gamma^{td} = 1 \iff N \mid td \iff
(N/d) \mid t$. Таких $t \in \{0, \dots, N - 1\}$ ровно $d$, и это степени $\gamma^{N/d}$,
т.\,е. $H_d = \langle \gamma^{N/d} \rangle$ "--- подгруппа порядка $d$. Если $H$ "--- любая
подгруппа порядка $d$, то по теореме Лагранжа $y^d = 1$ для всех $y \in H$, т.\,е.
$H \subseteq H_d$, а так как $|H| = |H_d|$, то $H = H_d$. Элемент $g$ порядка $d$ порождает
подгруппу $\langle g \rangle$ порядка $d$, т.\,е. $\langle g \rangle = H_d$.
\end{proof}

\checkyourself{Сколько элементов порядка $d$ в циклической группе порядка $N$, если $d \mid N$?
Почему условие $h^n = 1$ в предложении~\ref{l1:subgroups} не гарантирует $h \in \langle g
\rangle$ в нециклической группе (например, в $\Znx{8}$)?}

\section{Задача дискретного логарифмирования}

\begin{definition}
Пусть $g \in G$, $h \in G$. \emph{Задача дискретного логарифмирования} (DLP) с основанием $g$
состоит в нахождении такого целого $x$, что
\[
  g^x = h.
\]
Если решение существует, оно единственно по модулю $n = \ord(g)$ (предложение~\ref{l1:order}),
и наименьшее неотрицательное решение $x \in \{0, \dots, n - 1\}$ называется \emph{дискретным
логарифмом} $h$ по основанию $g$: $x = \log_g h$ (также $\ind_g h$ "--- \emph{индекс}).
\end{definition}

Классический случай "--- $G = \Znx{p}$: дано простое $p$ и вычеты $g, h$, найти $x$ с
$\congr{g^x}{h}{p}$. Решение существует для всех $h$ тогда и только тогда, когда $g$ "---
первообразный корень; иначе по предложению~\ref{l1:subgroups} оно есть ровно для $h$ с
$h^{\ord g} = 1$. Более общо, рассматривают DLP в $\Znx{n}$ для составного $n$, в мультипликативной
группе конечного поля $\GF{q}^*$, в группе точек эллиптической кривой (лекция~9) и т.\,д.

\begin{example}\label{l1:ex13}
$p = 13$, $g = 2$. Степени $2^k \bmod 13$:
\[
\begin{array}{c|cccccccccccc}
  k & 0 & 1 & 2 & 3 & 4 & 5 & 6 & 7 & 8 & 9 & 10 & 11 \\ \hline
  2^k & 1 & 2 & 4 & 8 & 3 & 6 & 12 & 11 & 9 & 5 & 10 & 7
\end{array}
\]
Все $12$ вычетов встретились, значит, $2$ "--- первообразный корень по модулю $13$, и, например,
$\log_2 7 = 11$, $\log_2 3 = 4$, $\log_2 5 = 9$. А уравнение $3^x \equiv 2 \pmod{13}$ решений
не имеет: $\ord(3) = 3$, $\langle 3 \rangle = \{1, 3, 9\} \not\ni 2$.
\end{example}

\begin{proposition}[свойства логарифма]\label{l1:loghom}
Пусть $\ord(g) = n$. Отображение
\[
  \log_g \colon \langle g \rangle \to \Zn{n}, \qquad h \mapsto \log_g h \bmod n,
\]
"--- изоморфизм групп (мультипликативной $\langle g \rangle$ и аддитивной $\Zn{n}$), т.\,е.
\[
  \log_g(h_1 h_2) \equiv \log_g h_1 + \log_g h_2, \qquad \log_g(h^k) \equiv k\log_g h,
  \qquad \log_g 1 = 0 \pmod n.
\]
Если $g'$ "--- другой образующий $\langle g \rangle$, то $\log_{g'} h \equiv \log_g h \cdot
(\log_g g')^{-1} \pmod n$.
\end{proposition}

\begin{proof}
Отображение $\exp_g \colon \Zn{n} \to \langle g \rangle$, $k \mapsto g^k$, корректно и
инъективно по предложению~\ref{l1:order}, сюръективно по определению $\langle g \rangle$ и
переводит сложение в умножение: $g^{k + l} = g^k g^l$. Значит, $\exp_g$ "--- изоморфизм, и
обратное к нему отображение $\log_g$ "--- тоже. Формулы для произведения и степени "--- запись
того, что $\log_g$ гомоморфизм. Наконец, если $g' = g^t$ "--- образующий, то $\gcd(t, n) = 1$
(иначе $\ord(g^t) = n/\gcd(t, n) < n$), и из $h = g'^{y} = g^{ty}$ получаем $\log_g h \equiv
ty$, откуда $y \equiv \log_g h \cdot t^{-1} \pmod n$.
\end{proof}

В примере~\ref{l1:ex13}: $\log_2 15 = \log_2 2 = 1$ и $\log_2 3 + \log_2 5 = 4 + 9 = 13 \equiv 1
\pmod{12}$. Обратите внимание: логарифмы складываются по модулю $p - 1$, а не по модулю $p$.

\begin{remark}
Сложность DLP зависит не от абстрактной группы, а от её \emph{представления}. Все циклические
группы порядка $n$ изоморфны $(\Zn{n}, +)$, но в $(\Zn{n}, +)$ «дискретный логарифм» "---
решение уравнения $xg \equiv h \pmod n$ "--- находится расширенным алгоритмом Евклида за
полиномиальное время. Трудность DLP в $\Znx{p}$ или на эллиптической кривой означает, что
изоморфизм $\log_g$ трудно \emph{вычислять}.
\end{remark}

\checkyourself{Почему $\log_g$ нельзя рассматривать как отображение в $\Zn{p}$? Чему равен
$\log_g(-1)$ в $\Znx{p}$, если $g$ "--- первообразный корень и $p > 2$?}

\section{Протокол Диффи"--~Хеллмана}

\begin{algorithm}[протокол Диффи"--~Хеллмана]\label{l1:dh}
\begin{algorithmic}[1]
\Require открытые параметры: простое $p$ и $g \in \Znx{p}$ порядка $n$ (обычно $n$ "--- большое
простое число).
\Ensure общий секрет $K$ Алисы и Боба.
\State Алиса выбирает секретное $a \in [1, n - 1]$ и отправляет Бобу $A = g^a \bmod p$.
\State Боб выбирает секретное $b \in [1, n - 1]$ и отправляет Алисе $B = g^b \bmod p$.
\State Алиса вычисляет $K = B^a \bmod p$, Боб вычисляет $K = A^b \bmod p$.
\end{algorithmic}
\end{algorithm}

\begin{figure}[h]
\centering
\begin{tikzpicture}[x=1cm,y=1cm]
  \node at (5,1.5) {открыто: $p$, $g$};
  \node at (0,0) {\personicon{caBlue}{Алиса}};
  \node at (10,0) {\personicon{caRed}{Боб}};
  \draw[-{Stealth[length=3mm]},thick] (1,0.6) -- (9,0.6) node[midway,above] {$A = g^a$};
  \draw[{Stealth[length=3mm]}-,thick] (1,0.0) -- (9,0.0) node[midway,below] {$B = g^b$};
  \node[font=\small,align=left,anchor=north] at (0,-0.8) {секрет $a$\\ $K = B^a = g^{ab}$};
  \node[font=\small,align=left,anchor=north] at (10,-0.8) {секрет $b$\\ $K = A^b = g^{ab}$};
  \node at (5,-2.2) {\personicon{caGray}{Ева}};
  \draw[dashed,caGray] (5,-0.6) -- (5,-1.4);
  \node[font=\small,anchor=north] at (5,-2.7) {знает $p, g, A, B$};
\end{tikzpicture}
\caption{Согласование ключа Диффи"--~Хеллмана}
\end{figure}

\begin{proposition}
Алиса и Боб получают один и тот же ключ $K = g^{ab}$, затратив по два возведения в степень,
т.\,е. $\bigO(\log n)$ умножений по модулю $p$ каждый.
\end{proposition}

\begin{proof}
$B^a = (g^b)^a = g^{ab} = (g^a)^b = A^b$ в $\Znx{p}$; оценка числа умножений "--- оценка
бинарного алгоритма возведения в степень.
\end{proof}

\begin{example}[{\cite[разд.~2.3]{hps}}]
$p = 941$, $g = 627$ (первообразный корень, $p - 1 = 2^2 \cdot 5 \cdot 47$). Алиса выбирает
$a = 347$ и отправляет $A = 627^{347} \equiv 390$, Боб выбирает $b = 781$ и отправляет
$B = 627^{781} \equiv 691$. Общий ключ $K = 691^{347} \equiv 390^{781} \equiv 470 \pmod{941}$.
\end{example}

\taskexample{dh-toy}

\subsection{Что знает Ева: задачи CDH и DDH}

Ева видит $p, g, A = g^a, B = g^b$ и хочет найти $K = g^{ab}$.

\begin{definition}
\begin{enumerate}
  \item \emph{Вычислительная задача Диффи"--~Хеллмана} (CDH): по $g, g^a, g^b$ найти $g^{ab}$.
  \item \emph{Задача распознавания Диффи"--~Хеллмана} (DDH): по $g, g^a, g^b, g^c$ определить,
    верно ли, что $g^c = g^{ab}$.
\end{enumerate}
\end{definition}

\begin{proposition}\label{l1:reductions}
DDH $\le$ CDH $\le$ DLP: алгоритм для DLP даёт алгоритм для CDH, а алгоритм для CDH "--- для DDH
(с дополнительными затратами в одно возведение в степень и одно сравнение).
\end{proposition}

\begin{proof}
Если мы умеем находить логарифмы, то по $A = g^a$ находим $a' = \log_g A \equiv a \pmod n$
и вычисляем $B^{a'} = g^{ab}$, т.\,е. решаем CDH. Если мы умеем решать CDH, то по $(g^a, g^b)$
находим $g^{ab}$ и сравниваем с данным $g^c$, т.\,е. решаем DDH.
\end{proof}

Обратные сведения в общем случае неизвестны: эквивалентность CDH и DLP доказана лишь для
некоторых классов групп (У.~Маурер и С.~Вольф, см.~\cite[гл.~21]{galbraith}), а DDH бывает
\emph{строго} проще, как показывает следующее утверждение.

\begin{proposition}\label{l1:ddh-easy}
Если $g$ "--- первообразный корень по модулю простого $p > 2$, то по $g^a, g^b, g^c$ можно за
$\bigO(\log p)$ умножений узнать чётности $a, b, c$; в частности, при нечётном $ab$ и чётном $c$
(и наоборот) DDH в $\Znx{p}$ решается сразу.
\end{proposition}

\begin{proof}
Элемент $\varepsilon = g^{(p-1)/2}$ удовлетворяет $\varepsilon^2 = 1$, $\varepsilon \ne 1$, а в поле
$\Zn{p}$ уравнение $y^2 = 1$ имеет лишь корни $\pm 1$, так что $\varepsilon = -1$. Поэтому
$(g^a)^{(p-1)/2} = (-1)^a$ определяет чётность $a$ (и аналогично для $b$, $c$), а чётность $ab$
равна чётности $c$, если $g^c = g^{ab}$ (так как $p - 1$ чётно).
\end{proof}

Поэтому на практике работают не во всей группе $\Znx{p}$, а в её подгруппе \emph{большого
простого порядка} $q$, где такой «утечки» нет (подробнее "--- в разделе~\ref{l1:params}).

\begin{remark}
Протокол в таком виде защищает только от \emph{пассивного} противника. Активная Ева может
подменить $A$ и $B$ своими $g^{e}$ (атака «человек посередине») и договориться с каждым из
участников отдельно. Поэтому в реальных протоколах (TLS, SSH) сообщения DH подписываются.
\end{remark}

\checkyourself{Алиса по ошибке выбрала $a = \ord(g)$. Что увидит Ева и каким будет общий ключ?
А если $g$ имеет порядок $5$?}

\section{Шифрование Эль-Гамаля}

Протокол Диффи"--~Хеллмана лишь согласует ключ. Т.~Эль-Гамаль~\cite{elgamal} (1985) построил на
той же идее систему шифрования с открытым ключом.

\begin{algorithm}[шифрование Эль-Гамаля]\label{l1:elgamal}
\begin{algorithmic}[1]
\Require открытые параметры $p$ и $g \in \Znx{p}$ порядка $n$.
\State \textbf{Генерация ключей} (Алиса): секретный ключ $a \in [1, n - 1]$, открытый ключ
$A = g^a \bmod p$.
\State \textbf{Шифрование} (Боб, сообщение $m \in \Znx{p}$): выбрать случайное одноразовое
$k \in [1, n - 1]$, вычислить $c_1 = g^k$, $c_2 = m A^k$; шифртекст "--- $(c_1, c_2)$.
\State \textbf{Расшифрование} (Алиса): $m = c_2 \cdot (c_1^{a})^{-1} \bmod p$.
\end{algorithmic}
\end{algorithm}

\begin{proposition}
Расшифрование корректно, а обратный элемент можно вычислить без алгоритма Евклида:
$(c_1^a)^{-1} = c_1^{p - 1 - a}$.
\end{proposition}

\begin{proof}
$c_1^a = g^{ka} = (g^a)^k = A^k$, поэтому $c_2 (c_1^a)^{-1} = m A^k A^{-k} = m$. По малой
теореме Ферма $c_1^{p-1} = 1$, откуда $c_1^{p-1-a} \cdot c_1^a = 1$.
\end{proof}

Иначе говоря, Боб проводит с Алисой обмен Диффи"--~Хеллмана с одноразовым секретом $k$ и маскирует
сообщение общим ключом $s = A^k = g^{ak}$.

\begin{example}[{\cite[разд.~2.4]{hps}}]
$p = 467$, $g = 2$, $a = 153$, $A = 2^{153} \equiv 224$. Боб шифрует $m = 331$ с $k = 197$:
$c_1 = 2^{197} \equiv 87$, $c_2 = 331 \cdot 224^{197} \equiv 57$. Алиса вычисляет $c_1^a =
87^{153} \equiv 367$, $367^{-1} \equiv 14$ и $m = 57 \cdot 14 \equiv 331 \pmod{467}$.
\end{example}

\taskexample{elgamal-toy}

\begin{proposition}[{\cite[разд.~2.4]{hps}}]
Расшифровать любой шифртекст Эль-Гамаля без секретного ключа можно тогда и только тогда, когда
можно решать CDH для пар $(A, c_1)$.
\end{proposition}

\begin{proof}
Если мы умеем решать CDH, то по $A = g^a$, $c_1 = g^k$ найдём $s = g^{ak}$ и $m = c_2 s^{-1}$.
Обратно, пусть есть «оракул», расшифровывающий любые пары $(c_1, c_2)$ для открытого ключа $A$.
Чтобы найти $g^{ab}$ по $A = g^a$ и $B = g^b$, подадим ему шифртекст $(B, 1)$: он вернёт
$m = 1 \cdot (B^a)^{-1}$, откуда $g^{ab} = m^{-1}$.
\end{proof}

\begin{remark}
Одноразовый показатель $k$ нельзя использовать повторно: если $(c_1, c_2)$ и $(c_1, c_2')$
зашифрованы с одним $k$ и Ева знает $m$, то $m' = m\, c_2' c_2^{-1}$. Шифртекст вдвое длиннее
сообщения "--- плата за вероятностное шифрование: одно и то же $m$ при разных $k$ шифруется
по-разному.
\end{remark}

\checkyourself{Что произойдёт, если Боб выберет $k$, кратное $\ord(g)$? Чем опасна маленькая
подгруппа $\langle g \rangle$ для шифрования Эль-Гамаля?}

\section{Тривиальная оценка сложности DLP}

Сложность алгоритмов DLP измеряют числом \emph{групповых операций} (умножений по модулю $p$);
каждая из них стоит $\bigO(\log^2 p)$ битовых операций школьным методом (быстрее "--- методами
лекций 2 и 4 первого семестра). Размер входа "--- $\log_2 p$ бит, так что «полиномиальный»
алгоритм должен делать $\bigO(\log^c p)$ операций.

\begin{algorithm}[перебор]\label{l1:brute}
\begin{algorithmic}[1]
\Require $g, h \in G$, $n = \ord(g)$ (или верхняя граница для него, например $|G|$).
\Ensure $x = \log_g h$ или «решений нет».
\State $y \gets 1$
\For{$x = 0, 1, \dots, n - 1$}
  \If{$y = h$} \Return $x$ \EndIf
  \State $y \gets y \cdot g$
\EndFor
\State \Return «решений нет»
\end{algorithmic}
\end{algorithm}

\begin{proposition}
Перебор корректен и делает не более $n - 1$ умножений и $n$ сравнений, храня $\bigO(1)$
элементов группы. Для случайного $h \in \langle g \rangle$ среднее число умножений "--- около
$n/2$.
\end{proposition}

\begin{proof}
На шаге с номером $x$ выполнено $y = g^x$, и перебираются все $x \in [0, n - 1]$, т.\,е. все
элементы $\langle g \rangle$ (предложение~\ref{l1:order}). Если $h = g^{x_0}$ с равномерно
распределённым $x_0 \in [0, n - 1]$, то число умножений равно $x_0$, а $\mathsf{E}\,x_0 =
(n - 1)/2$.
\end{proof}

Для $n \approx p \approx 2^{2048}$ это $2^{2047}$ операций "--- экспонента от длины входа.
Заметим, что перебор попутно находит $\ord(g)$ (первое $x > 0$ с $g^x = 1$) и все решения
$x \equiv x_0 \pmod n$.

\taskexample{dlp-brute-force}

\section{Выбор параметров}\label{l1:params}

Будущие лекции покажут, что сложность DLP определяется не размером группы, а её «худшей частью»:
\begin{itemize}
  \item алгоритм Полига"--~Хеллмана (лекция~2) сводит DLP в группе порядка $n =
    \prod q_i^{e_i}$ к задачам в подгруппах простых порядков $q_i$, поэтому время определяется
    \emph{наибольшим} простым делителем $q$ числа $\ord(g)$;
  \item универсальные алгоритмы (Шенкс, Поллард) решают DLP в группе простого порядка $q$ за
    $\bigO(\sqrt q)$ операций, и это неулучшаемо для универсальных алгоритмов (раздел~\ref{l1:generic});
  \item для $\Znx{p}$ есть субэкспоненциальное исчисление индексов (лекция~4), поэтому сам модуль
    $p$ должен быть значительно длиннее, чем $q$.
\end{itemize}
Отсюда стандартный выбор: $p$ длиной $2048$"--~$3072$ бит и $g$ порядка $q$, где $q$ "---
простое длиной $256$ бит (или больше), $q \mid p - 1$. Частный случай "--- \emph{безопасное
простое} $p = 2q + 1$ с простым $q$.

\begin{proposition}\label{l1:safe}
Пусть $p = 2q + 1$, $p$ и $q$ "--- простые, $q > 2$. Тогда порядки элементов $\Znx{p}$ "--- это
$1, 2, q, 2q$; элемент $y \ne 1$ имеет порядок $q$ тогда и только тогда, когда он является
квадратом. В частности, $\ord(4) = q$.
\end{proposition}

\begin{proof}
Порядки делят $p - 1 = 2q$. Квадраты $y = z^2$ удовлетворяют $y^q = z^{2q} = z^{p-1} = 1$, т.\,е.
лежат в подгруппе $H_q$ порядка $q$ (предложение~\ref{l1:subgroups}); квадратов в $\Znx{p}$ ровно
$(p-1)/2 = q$ (отображение $z \mapsto z^2$ склеивает $z$ и $-z$), так что $H_q$ "--- в точности
квадраты. Если $y \in H_q$, $y \ne 1$, то $\ord(y) \mid q$ и $\ord(y) \ne 1$, т.\,е. $\ord(y) = q$.
Наконец, $4 = 2^2 \ne 1$ при $p > 3$.
\end{proof}

Например, $p = 23 = 2 \cdot 11 + 1$: $\ord(4) = 11$, а $\ord(22) = 2$ ($22 \equiv -1$). Получив от
партнёра $A$, участник проверяет $1 < A < p - 1$ и $A^q \equiv 1$; иначе противник может подсунуть
элемент маленького порядка, и общий ключ примет лишь несколько значений (атака «заключения
в малую подгруппу»).

\begin{remark}[уровни стойкости]
Рекомендации NIST~\cite[табл.~2]{nist-sp800-57} сопоставляют параметры с длиной ключа
симметричного шифра той же стойкости ($\log_2$ числа операций лучшей известной атаки):
\begin{center}
\begin{tabular}{c|c|c|c}
  стойкость, бит & DH, Эль-Гамаль: $\log_2 p$ & $\log_2 q$ & эллиптические кривые: $\log_2 q$ \\ \hline
  $80$  & $1024$  & $160$ & $160$ \\
  $112$ & $2048$  & $224$ & $224$ \\
  $128$ & $3072$  & $256$ & $256$ \\
  $192$ & $7680$  & $384$ & $384$ \\
  $256$ & $15360$ & $512$ & $512$
\end{tabular}
\end{center}
Длина $q$ "--- это удвоенная стойкость (из-за алгоритмов за $\sqrt q$), а модуль $p$ гораздо
длиннее из-за исчисления индексов. Стойкость $80$ бит сегодня считается недостаточной.
\end{remark}

\checkyourself{Почему для безопасного простого $p = 2q + 1$ нельзя брать $g = p - 1$? Сколько
значений может принять общий ключ Диффи"--~Хеллмана, если $g = -1$?}

\section{Алгоритм Шенкса (baby-step giant-step)}

Идея Д.~Шенкса~\cite{shanks} "--- обмен памяти на время: составить два списка длины около
$\sqrt n$ так, чтобы их общий элемент давал логарифм.

\begin{lemma}[о коллизии]\label{l1:collision}
Пусть $\ord(g) = n$, $m \ge 1$ "--- целое с $m^2 \ge n$ (например, $m = \ceil{\sqrt n}$), и
\[
\begin{gathered}
  L_1 = \{g^j : 0 \le j < m\} \quad\text{(«малые шаги»)},\\
  L_2 = \{h \cdot g^{-im} : 0 \le i < m\} \quad\text{(«большие шаги»)}.
\end{gathered}
\]
Тогда $L_1 \cap L_2 \ne \varnothing$ $\iff$ $h \in \langle g \rangle$. При этом если
$g^j = h g^{-im}$, то $x = im + j$ "--- решение уравнения $g^x = h$.
\end{lemma}

\begin{proof}
Если $g^j = h g^{-im}$, то $h = g^{im + j} \in \langle g \rangle$ и $x = im + j$ "--- решение.
Обратно, пусть $h = g^x$, где $0 \le x < n$ (предложение~\ref{l1:order}). Разделим $x$ на $m$ с
остатком: $x = im + j$, $0 \le j < m$. Тогда $i = \floor{x/m} \le (n - 1)/m < m^2/m = m$, т.\,е.
$0 \le i < m$, и $h g^{-im} = g^{x - im} = g^j$ "--- общий элемент $L_1$ и $L_2$.
\end{proof}

\begin{algorithm}[Шенкс]\label{l1:bsgs}
\begin{algorithmic}[1]
\Require $g, h \in G$; $n = \ord(g)$ или верхняя граница для него (например, $p - 1$).
\Ensure $x$ с $g^x = h$ или «решений нет».
\State $m \gets \ceil{\sqrt n}$
\State $T \gets$ пустая хеш-таблица; $y \gets 1$
\For{$j = 0, 1, \dots, m - 1$} \Comment{малые шаги}
  \State если ключа $y$ нет в $T$, то $T[y] \gets j$; \quad $y \gets y \cdot g$
\EndFor
\State $u \gets g^{-m}$ \Comment{$g^{-1}$ "--- алгоритмом Евклида или как $g^{n-1}$}
\State $y \gets h$
\For{$i = 0, 1, \dots, m - 1$} \Comment{большие шаги}
  \If{$y \in T$} \Return $im + T[y]$ \EndIf
  \State $y \gets y \cdot u$
\EndFor
\State \Return «решений нет»
\end{algorithmic}
\end{algorithm}

\begin{theorem}\label{l1:bsgs-complexity}
Алгоритм Шенкса корректен и находит дискретный логарифм за $\bigO(\sqrt n)$ групповых
операций (не более $2m + \bigO(\log n)$ умножений), используя память на $m = \bigO(\sqrt n)$
элементов группы и $\bigO(\sqrt n)$ обращений к хеш-таблице. Если вместо хеш-таблицы список $L_1$
сортировать, а элементы $L_2$ искать двоичным поиском, время составит
$\bigO(\sqrt n \log n)$ сравнений.
\end{theorem}

\begin{proof}
Корректность "--- лемма~\ref{l1:collision}: таблица $T$ содержит $L_1$ (с наименьшим $j$ для
каждого значения), а во втором цикле перебираются элементы $L_2$ в порядке возрастания $i$;
верхняя граница $N \ge n$ вместо $n$ тоже подходит, так как $m^2 \ge N \ge n$. Малые шаги
требуют $m - 1$ умножений, вычисление $u = g^{-m}$ "--- $\bigO(\log m)$ умножений и одно обращение,
большие шаги "--- не более $m - 1$ умножений. Хранится $m$ элементов. Сортировка $m$ элементов
стоит $\bigO(m \log m)$ сравнений, $m$ двоичных поисков "--- столько же.
\end{proof}

Битовая сложность при $G = \Znx{p}$ и $n \approx p$: $\bigO(\sqrt p \log^2 p)$ "--- всё ещё
экспонента от длины $\log p$, но с вдвое меньшим показателем, чем у перебора. Для $p \approx 2^{2048}$
это $2^{1024}$ операций (недостижимо), но для группы порядка $2^{80}$ "--- лишь $2^{40}$
(несколько минут на ноутбуке), если хватает памяти на $2^{40}$ элементов.

\begin{remark}[варианты]
\begin{enumerate}
  \item Длины списков можно взять разными: при $m_1 m_2 \ge n$ подходят $L_1 = \{g^j : j < m_1\}$
    и $L_2 = \{h g^{-i m_1} : i < m_2\}$. Время $m_1 + m_2$ при фиксированном произведении
    минимально при $m_1 = m_2 = \sqrt n$; уменьшая $m_1$, экономим память ценой времени.
  \item Можно хранить большие шаги $g^{im}$, $1 \le i \le m$, и искать в них $h g^{j}$ "---
    тогда $x = im - j$; это вариант \cite[разд.~2.7]{hps}.
  \item Если известно, что $x$ лежит в интервале длины $w$ (например, показатель «короткий»),
    достаточно $m = \ceil{\sqrt w}$: сложность $\bigO(\sqrt w)$ вместо $\bigO(\sqrt n)$.
  \item Если $\ord(g)$ известен и меньше $p - 1$, длину списков выбирают по $\ord(g)$:
    $\ceil{\sqrt{\ord g}}$ вместо $\ceil{\sqrt{p - 1}}$.
\end{enumerate}
\end{remark}

\begin{example}
$p = 101$, $g = 2$ (первообразный корень, $n = 100$), решим $2^x \equiv 81 \pmod{101}$.
$m = 10$, малые шаги $2^j$:
\[
\begin{array}{c|cccccccccc}
  j & 0 & 1 & 2 & 3 & 4 & 5 & 6 & 7 & 8 & 9 \\ \hline
  2^j & 1 & 2 & 4 & 8 & 16 & 32 & 64 & 27 & 54 & 7
\end{array}
\]
$2^{10} \equiv 14$, $g^{-10} \equiv 14^{-1} \equiv 65$ (так как $14 \cdot 65 = 910 = 9 \cdot 101 +
1$). Большие шаги $y_i = 81 \cdot 65^i$:
\[
\begin{array}{c|cccccccc}
  i & 0 & 1 & 2 & 3 & 4 & 5 & 6 & 7 \\ \hline
  y_i & 81 & 13 & 37 & 82 & 78 & 20 & 88 & \mathbf{64}
\end{array}
\]
$y_7 = 64 = 2^6$, откуда $x = 7 \cdot 10 + 6 = 76$. Перебору понадобилось бы $76$ умножений,
здесь "--- $9 + 7$ плюс вычисление $g^{-10}$.
\end{example}

\taskexample{bsgs-shanks}

\taskexample{bsgs-subgroup}

\checkyourself{Почему в лемме~\ref{l1:collision} достаточно $m^2 \ge n$, и что сломается при
$m^2 < n$? Сколько памяти понадобится алгоритму Шенкса для группы порядка $2^{128}$?}

\checkyourself{Как изменить алгоритм Шенкса, если $\ord(g)$ неизвестен, а известна лишь граница
$\ord(g) \le N$? Что вернёт алгоритм, если $h \notin \langle g \rangle$?}

\section{Нижняя оценка для универсальных алгоритмов}\label{l1:generic}

Перебор и алгоритм Шенкса \emph{универсальны} (generic): они используют только групповые
операции и сравнение элементов и не «заглядывают» в их представление. В модели \emph{чёрного
ящика} группа $\Zn{n}$ задаётся случайным инъективным кодированием $\sigma \colon \Zn{n} \to
\bits{\ell}$, а алгоритм получает коды элементов и может обращаться к оракулу, который по
$\sigma(x), \sigma(y)$ возвращает $\sigma(x \pm y)$.

\begin{theorem}[В.~Шуп~\cite{shoup-generic}, без доказательства]\label{l1:shoup}
Пусть $q$ "--- наибольший простой делитель $n$. Любой универсальный (вероятностный) алгоритм,
решающий DLP в циклической группе порядка $n$ с вероятностью не меньше $1/2$, выполняет
$\Omega(\sqrt q)$ групповых операций.
\end{theorem}

Таким образом, в группе простого порядка алгоритм Шенкса (как и $\rho$-метод Полларда,
лекция~3, которому нужна лишь $\bigO(1)$ память) оптимален среди универсальных алгоритмов, а
в группе составного порядка оценка $\sqrt q$ достигается алгоритмом Полига"--~Хеллмана
(лекция~2). Обойти оценку можно, только используя представление элементов: для $\Znx{p}$ "---
разложение чисел на простые множители (исчисление индексов, лекция~4). Для групп точек
«хороших» эллиптических кривых подобных методов не известно, поэтому там хватает $q \approx 2^{256}$
при той же стойкости, что и $p \approx 2^{3072}$ в $\Znx{p}$.

\checkyourself{Противоречит ли теорема~\ref{l1:shoup} существованию субэкспоненциального
алгоритма исчисления индексов для $\Znx{p}$? Почему при оценке сложности DLP в группе порядка
$n = 2^{100} \cdot 3$ важен не размер $n$, а его простые делители?}

\section{Итоги}

\begin{itemize}
  \item DLP: по $g, h$ найти $x$ с $g^x = h$; решение единственно по модулю $\ord(g)$ и
    существует ровно для $h \in \langle g \rangle$, т.\,е. при $h^{\ord g} = 1$.
  \item Протокол Диффи"--~Хеллмана и шифрование Эль-Гамаля стойки не более, чем DLP; их стойкость
    "--- это трудность CDH.
  \item Перебор: $\bigO(n)$ операций и $\bigO(1)$ памяти; алгоритм Шенкса: $\bigO(\sqrt n)$
    операций и $\bigO(\sqrt n)$ памяти; для универсальных алгоритмов $\Omega(\sqrt q)$ "---
    нижняя оценка.
  \item Параметры: подгруппа большого простого порядка $q$; проверка принадлежности $A^q = 1$;
    $\log_2 q \approx 2 \times$ (стойкость в битах).
\end{itemize}

\checkyourselfsection

\begin{furtherreading}
  \readitem[разд.~2.2--2.4, 2.6, 2.7]{hps}{DLP, Диффи"--~Хеллман, Эль-Гамаль, оценки сложности,
    алгоритм Шенкса (основной источник лекции)}
  \readitem[разд.~13.1, 13.3; гл.~20, 21]{galbraith}{перебор и baby-step giant-step с вариантами,
    задачи CDH и DDH, сведения между ними}
  \readitem[разд.~3.6, 8.4, 12.6]{hac}{алгоритмы DLP, шифрование Эль-Гамаля, протокол
    Диффи"--~Хеллмана}
  \readitem[гл.~8]{paar-pelzl}{протоколы на основе DLP и выбор параметров на практике}
  \readitem[разд.~11.2]{shoup}{алгоритм Шенкса в $\Znx{p}$ и его анализ}
  \readitem{shoup-generic}{нижняя оценка для универсальных алгоритмов}
  \readitem[табл.~2]{nist-sp800-57}{соответствие длин ключей и уровней стойкости}
  \readitem{diffie-hellman}{оригинальная статья Диффи и Хеллмана}
  \readitem[неделя~5]{coursera-boneh}{видеолекции: согласование ключа, Диффи"--~Хеллман}
  \readitem[разд.~«Discrete Logarithm»]{cp-algorithms}{реализация алгоритма Шенкса с хеш-таблицей}
\end{furtherreading}
\end{lecture}
